% !PRINTED 591
\weberpage{591}
\begin{center}
\S.\,161.\\[4pt]
\gesperrt{Primitivwurzeln der Primideale.}
\end{center}

Aus den Formeln (4), (5) des vorigen Paragraphen ergiebt sich, mit Rücksicht auf \S.\,150, für jede ganze Zahl $\eta$ des Körpers $R$
\[
\eta^P \equiv \eta \pmod{\mathfrak{P}},
\]
also auch
\weq{1}{\eta^P \equiv \eta \pmod{\mathfrak{p}},}
und für jede Zahl $\omega$ in $\mathfrak{o}$
\weq{2}{\omega^{P^f} \equiv \omega \pmod{\mathfrak{p}}.}
Die Anzahl der nach $\mathfrak{P}$ incongruenten ganzen Zahlen in $R$ ist $P$, und die Anzahl der nach $\mathfrak{p}$ incongruenten Zahlen in $\Omega$ ist $P^f$.

Wir haben schon früher (\S.\,150) gezeigt, dass es zu jedem Primideal $\mathfrak{p}$ \gesperrt{Primitivwurzeln} $\gamma$ giebt, d. h. Zahlen in $\Omega$, die der Congruenz
\weq{3}{\gamma^{P^f-1} \equiv 1 \pmod{\mathfrak{p}}}
genügen, und zugleich die Eigenschaft haben, dass keine niedrigere Potenz mit positiven Exponenten der Einheit congruent wird. Dann ist jede durch $\mathfrak{p}$ nicht theilbare ganze Zahl in $\Omega$ mit einer und nur mit einer der Potenzen von $\gamma$
\weq{4}{1, \gamma, \gamma^2, \ldots, \gamma^{P^f-2}}
nach dem Modul $\mathfrak{p}$ congruent.

Jede ganze Zahl $\Theta$ in $\Omega$ genügt einer Gleichung höchstens vom $n^{\text{ten}}$ Grade, deren Coëfficienten ganze Zahlen in $R$ sind.

Diese Gleichung ist zugleich eine Congruenz nach dem Modul $\mathfrak{p}$, und unter allen solchen Congruenzen wird es eine von \gesperrt{möglichst niedrigem Grade}
\[
F(\Theta) \equiv 0 \pmod{\mathfrak{p}}
\]
geben, in der wir überdies den Coëfficienten der höchsten Potenz von $\Theta$ gleich 1 annehmen können. Denn ist der höchste (durch $\mathfrak{p}$ und folglich durch $\mathfrak{P}$ untheilbare) Coëfficient $\eta_0$, so können wir immer der Congruenz
\[
\eta_0 \eta \equiv 1 \pmod{\mathfrak{P}}
\]
durch ein ganzzahliges $\eta$ in $R$ genügen, und haben dann nur die Function $F$ mit diesem $\eta$ zu multipliciren. Wir können
%CONT>>
